\documentclass[11pt,a4paper]{article}
\usepackage[top=1.7cm,bottom=1.5cm,left=1.5cm,right=1.5cm,headsep=8pt,footskip=16pt]{geometry}
\usepackage{enumitem}
\usepackage{amssymb}
\usepackage{array}
\usepackage[table]{xcolor}
\usepackage{fancyhdr}
\usepackage{lastpage}
\usepackage{multicol}
\usepackage{polyglossia}
\setdefaultlanguage[calendar=gregorian, numerals=maghrib]{arabic}
\setotherlanguage{english}
\IfFontExistsTF{Amiri}{%
  \newfontfamily\arabicfont[Script=Arabic, Scale=1.08]{Amiri}
}{%
  \newfontfamily\arabicfont[Script=Arabic, Scale=1.08, Path=../fonts/, Extension=.ttf,
    UprightFont=*-Regular, BoldFont=*-Bold, ItalicFont=*-Italic,
    BoldItalicFont=*-BoldItalic]{Amiri}
}
\IfFontExistsTF{AlBattar}{%
  \newfontfamily\khatt[Script=Arabic, Scale=1.15]{AlBattar}
}{%
  \newfontfamily\khatt[Script=Arabic, Scale=1.15, Path=../fonts/,
    BoldFont=ae_AlBattar.ttf]{ae_AlBattar.ttf}
}
\renewcommand{\headrulewidth}{0.8pt}
\renewcommand{\footrulewidth}{0.8pt}
\fancyhf{}
\rhead{\khatt المدرسة العليا للأساتذة ببوزريعة}
\lhead{\khatt ملحقة بومرداس}
\lfoot{\khatt الأستاذ: د. بن عبد الرزاق يوسف}
\cfoot{\thepage/\pageref{LastPage}}
\pagestyle{fancy}
\parindent 0pt
\setlength{\columnsep}{0.6cm}
\setlength{\columnseprule}{0.4pt}
\setlist[enumerate,1]{label=\textbf{\arabic*.}, leftmargin=*, itemsep=5pt, topsep=2pt}
\begin{document}

\linespread{1.08}\selectfont
\begin{center}
{\LARGE\khatt\textbf{الاختبار الثالث والعشرون في الذكاء الاصطناعي}}\\[4pt]
{\normalsize وضعيات إشكالية ــ أسئلة متعددة الاختيارات ــ الفصول من 1 إلى 6}
\end{center}
\vspace{-2pt}\rule{\textwidth}{0.6pt}\\[2pt]
\textbf{التعليمة:} يتضمن الاختبار خمس وضعيات إشكالية. اقرأ كل وضعية بعناية، ثم اختر التحليل أو التصرف الأنسب
من بين الاختيارات الخمسة؛ لكل وضعية إجابة صحيحة واحدة فقط.
اكتب \textbf{حرف} الإجابة الصحيحة في جدول الإجابات في آخر الاختبار.
لكل وضعية 4 نقاط (المجموع 20 نقطة).
\vspace{4pt}

\noindent
{\large\khatt\textbf{الوضعية 1}}\hfill{\small(الفصل 1)}\par\nopagebreak\vspace{3pt}
\setlength{\fboxsep}{7pt}\colorbox{gray!10}{\parbox{\dimexpr\linewidth-14pt}{طلب طالب من أداة ذكاء اصطناعي غير متصلة بالبحث: «لخّص لي آخر مقال نُشر في موقع مجلة جامعتنا»، وألصق رابطه فقط. فقدّمت الأداة ملخصًا مفصّلًا مقنعًا. وعندما قرأ المقال وجده يتناول موضوعًا مختلفًا تمامًا. ويتساءل كيف «قرأت» الأداة المقال. ويريد أن يستعمل الأداة مستقبلًا في تلخيص مقالات المجلة لزملائه. ما التفسير الأدق؟}}\par\nopagebreak
\begin{itemize}[leftmargin=1.0cm, labelsep=6pt, itemsep=3pt, topsep=4pt, parsep=0pt]
  \item[\textbf{أ)}] الأداة غير المتصلة لا تستطيع فتح الرابط، فولّدت ملخصًا محتملًا من كلمات الرابط وسياق الطلب؛ ولتلخيص مقال يجب لصق نصه أو رفعه، أو استعمال أداة متصلة بالبحث مع التحقق.
  \item[\textbf{ب)}] الأداة غير المتصلة لا تفتح الروابط، لكنها تعرف كل مقالات مجلة الجامعة من بيانات تدريبها، فالملخص يخص مقالًا قديمًا من المجلة نفسها، وهو صحيح بالنسبة إلى ذلك المقال.
  \item[\textbf{ج)}] الأداة قرأت المقال فعلًا، لكن الموقع غيّر محتواه بعد ذلك، فالملخص صحيح بالنسبة إلى المقال الذي كان منشورًا لحظة الطلب، ولا حاجة إلى لصق النص في المرات القادمة، كما يحدث كثيرًا في المواقع الإلكترونية.
  \item[\textbf{د)}] الأداة لا تستطيع تلخيص أي مقال إطلاقًا، حتى لو لُصق نصه كاملًا في الطلب، لأن التلخيص يتطلب فهمًا لا تملكه الأدوات، وكل الملخصات التي تنتجها مختلقة بالضرورة.
  \item[\textbf{هـ)}] الأداة قرأت مقالًا آخر في الموقع نفسه، لأن الروابط تتشابه، والحل إعادة لصق الرابط نفسه مرة ثانية، فتقرأ الأداة المقال الصحيح هذه المرة دون حاجة إلى لصق نصه.
\end{itemize}
\par\vspace{6pt}
\noindent
{\large\khatt\textbf{الوضعية 2}}\hfill{\small(الفصل 2)}\par\nopagebreak\vspace{3pt}
\setlength{\fboxsep}{7pt}\colorbox{gray!10}{\parbox{\dimexpr\linewidth-14pt}{درّب فريق نموذجًا يتنبأ بعلامة الإنشاء انطلاقًا من آلاف الإنشاءات المصحَّحة سابقًا. لاحظ أن النموذج يمنح علامات مرتفعة لكل نص طويل حتى لو كان مكررًا، وتبيّن أن الإنشاءات الطويلة في بيانات التدريب كانت غالبًا ذات علامات مرتفعة. ويريد الفريق تقديم النموذج للمؤسسات التربوية قريبًا. ما التفسير الأدق والإجراء الأنسب؟}}\par\nopagebreak
\begin{itemize}[leftmargin=1.0cm, labelsep=6pt, itemsep=3pt, topsep=4pt, parsep=0pt]
  \item[\textbf{أ)}] الخلل في الأساتذة الذين صححوا الإنشاءات، فهم منحازون إلى الطول، والحل تجاهل النموذج والاعتماد على تصحيحهم نفسه، لأن النموذج لا يمكن تصحيحه بأي تعديل في البيانات.
  \item[\textbf{ب)}] النموذج على صواب، لأن الطول سبب الجودة في الإنشاء دائمًا، والاقتران في البيانات دليل على هذه السببية، فلا حاجة إلى أي تعديل في البيانات أو في طريقة اختبار النموذج، ويبقى صالحًا للاستعمال في المؤسسات.
  \item[\textbf{ج)}] النموذج تعلّم اختصارًا إحصائيًا: ربط الطول بالعلامة لأنهما اقترنا في البيانات، لا لأن الطول سبب الجودة؛ فتُضاف أمثلة طويلة ضعيفة وقصيرة جيدة، ويُختبر النموذج على حالات تكسر هذا الاقتران.
  \item[\textbf{د)}] النموذج تعلّم اختصارًا إحصائيًا، والحل أن يُكتب له في التعليمات «لا تهتم بالطول»، لأن هذه التعليمة تلغي ما تعلّمه من البيانات، فلا حاجة إلى أمثلة جديدة أو اختبارات خاصة، فيصبح النموذج عادلًا مع كل الإنشاءات.
  \item[\textbf{هـ)}] النموذج تعلّم اختصارًا إحصائيًا، والحل حذف كل الإنشاءات الطويلة من البيانات، لأن غياب النصوص الطويلة يمنع النموذج من ربط الطول بالعلامة، دون حاجة إلى اختبار حالات خاصة.
\end{itemize}
\par\vspace{6pt}
\noindent
{\large\khatt\textbf{الوضعية 3}}\hfill{\small(الفصل 3)}\par\nopagebreak\vspace{3pt}
\setlength{\fboxsep}{7pt}\colorbox{gray!10}{\parbox{\dimexpr\linewidth-14pt}{يستعمل أستاذ أداة لاقتراح تواريخ ومعلومات عن شعراء مغمورين لدرس في الأدب الجزائري القديم. لاحظ أن الأداة تقدّم كل المعلومات بالدرجة نفسها من الثقة، فيصعب عليه تحديد ما يحتاج إلى تحقق دقيق، ووقته محدود. وهو يريد تقليل التحقق إلى الحد الأدنى دون أن يُدخل خطأ إلى درسه. أي تعديل في الطلب يساعده أكثر؟}}\par\nopagebreak
\begin{itemize}[leftmargin=1.0cm, labelsep=6pt, itemsep=3pt, topsep=4pt, parsep=0pt]
  \item[\textbf{أ)}] يحذف كل الشعراء المغمورين من الدرس، ويكتفي بالشعراء المشهورين الذين لا تخطئ الأداة في معلوماتهم أبدًا، لأن الشهرة تضمن صحة كل ما تقوله الأداة دون أي تحقق.
  \item[\textbf{ب)}] يطلب أن تميّز الأداة بين المعلومات المشهورة الموثقة والمعلومات التي قد تكون غير دقيقة، وأن تقترح مصادر للتحقق، ثم يتحقق من الجميع مع أولوية لما وُسم بعدم اليقين، فالتمييز يوجّه جهده ولا يغني عنه.
  \item[\textbf{ج)}] يطلب معلومات أكثر عن كل شاعر، لأن كثرة التفاصيل تدل على دقة المعلومات، فيعتمد الشعراء الذين قدّمت الأداة عنهم تفاصيل كثيرة، ويتحقق فقط ممن قدّمت عنهم تفاصيل قليلة، فيكون التحقق مركّزًا ومختصرًا.
  \item[\textbf{د)}] يطلب أن تميّز الأداة بين المعلومات الموثقة وغير الموثقة، ثم يتحقق فقط مما وسمته بعدم اليقين، ويعتمد الباقي مباشرة، لأن الأداة تعرف بدقة ما تعرفه وما لا تعرفه في كل الحالات، فيوفر وقتًا كبيرًا في التحضير.
  \item[\textbf{هـ)}] يطلب من الأداة أن تكون «واثقة دائمًا»، لأن الثقة الثابتة تجعل المعلومات متجانسة في الدرس، والتمييز بين درجات اليقين يربك التلاميذ ولا يفيد الأستاذ في التحضير.
\end{itemize}
\par\vspace{6pt}
\noindent
{\large\khatt\textbf{الوضعية 4}}\hfill{\small(الفصل 4)}\par\nopagebreak\vspace{3pt}
\setlength{\fboxsep}{7pt}\colorbox{gray!10}{\parbox{\dimexpr\linewidth-14pt}{اقترح عضو في إدارة مدرسة خاصة توليد «شهادات» لأولياء وتلاميذ وهميين بأداة ذكاء اصطناعي، مع صور وجوه مولَّدة، ونشرها في موقع المدرسة للترويج للتسجيلات. ويرى أن الأمر «تسويق عادي» ما دامت المدرسة جيدة فعلًا. ويعارضه عضو آخر يرى في ذلك خداعًا للأولياء، والمدير يطلب رأيًا مؤسسًا. ما التقدير الأخلاقي الأدق؟}}\par\nopagebreak
\begin{itemize}[leftmargin=1.0cm, labelsep=6pt, itemsep=3pt, topsep=4pt, parsep=0pt]
  \item[\textbf{أ)}] يجوز النشر إذا كُتب في أسفل الصفحة بخط صغير «بعض المحتوى قد يكون توضيحيًا»، لأن هذه العبارة العامة تكفي قانونيًا وأخلاقيًا، ولا يلزم التصريح بأن الشهادات مختلقة.
  \item[\textbf{ب)}] الإشكال الوحيد في الصور المولَّدة، فإذا حُذفت الصور جاز نشر الشهادات المختلقة بأسماء وهمية، لأن النص المكتوب لا يضلل الأولياء كما تضللهم صور الوجوه الواقعية.
  \item[\textbf{ج)}] يجوز النشر بشرط أن تكون الشهادات إيجابية معتدلة لا مبالغة فيها، لأن الإشكال الأخلاقي في المبالغة لا في الاختلاق، والاعتدال يجعل الشهادات المختلقة مقبولة في الترويج.
  \item[\textbf{د)}] نشر شهادات مختلقة على أنها حقيقية تضليل للأولياء ومخالفة لمبدأ الصدق والشفافية، ولو كانت المدرسة جيدة؛ والبديل جمع شهادات حقيقية بموافقة أصحابها، أو عرض معطيات موثقة.
  \item[\textbf{هـ)}] لا إشكال، لأن الشهادات تعبّر عن جودة حقيقية للمدرسة، والوجوه المولَّدة لا تمثل أشخاصًا حقيقيين، فلا يُنتهك حق أحد في صورته، والتسويق يقوم دائمًا على المبالغة المقبولة.
\end{itemize}
\par\vspace{6pt}
\noindent
{\large\khatt\textbf{الوضعية 5}}\hfill{\small(الفصل 5)}\par\nopagebreak\vspace{3pt}
\setlength{\fboxsep}{7pt}\colorbox{gray!10}{\parbox{\dimexpr\linewidth-14pt}{طلب أستاذ من أداة استخراج «الحقائق» من مقال رأي عن تعليم العربية في المدارس الخاصة، لبناء نشاط في التمييز بين الحقيقة والرأي. جاءت القائمة منظمة، لكنها تضمنت أحكامًا تقييمية للكاتب مقدَّمة على أنها حقائق، مثل «المناهج الحالية فاشلة». ويريد أن يقدّم النشاط في حصة الغد. ما التصرف الأنسب؟}}\par\nopagebreak
\begin{itemize}[leftmargin=1.0cm, labelsep=6pt, itemsep=3pt, topsep=4pt, parsep=0pt]
  \item[\textbf{أ)}] يعرض القائمة على التلاميذ على أنها نموذج صحيح للتمييز بين الحقيقة والرأي دون مراجعة، لأن أخطاء الأداة ستمر دون أن يلاحظها التلاميذ، فلا تؤثر في تعلّمهم.
  \item[\textbf{ب)}] يعتمد القائمة كما هي، لأن الكاتب متخصص، وأحكام المتخصصين حقائق بالضرورة، والتمييز بين الحقيقة والرأي لا ينطبق على كتابات الخبراء في مجال اختصاصهم.
  \item[\textbf{ج)}] يطلب من الأداة إعادة الاستخراج، ويعتمد القائمة الثانية دون مراجعة، لأن الأداة تميّز بدقة بين الحقيقة والرأي في المحاولة الثانية، فلا داعي لوضع معيار أو مراجعة الأحكام.
  \item[\textbf{د)}] يحذف من القائمة كل جملة تتضمن صفة، لأن الصفة دليل الرأي دائمًا، ويعتمد الباقي حقائق دون تحقق، لأن الجمل الخالية من الصفات حقائق قابلة للتحقق بالضرورة.
  \item[\textbf{هـ)}] يراجع القائمة بنفسه بمعيار واضح: الحقيقة قابلة للتحقق، والرأي حكم تقييمي، فيعيد تصنيف الأحكام المدرجة خطأً، ويمكنه استثمار أخطاء الأداة نفسها مادة للنقاش بعد تصحيحها.
\end{itemize}
\par\vspace{6pt}
\begin{center}
{\khatt\textbf{جدول الإجابات}}\\[3pt]
\renewcommand{\arraystretch}{1.6}
\begin{tabular}{|>{\centering\arraybackslash}m{1.9cm}|*{5}{>{\centering\arraybackslash}m{1.2cm}|}}
\hline
\textbf{رقم السؤال} & \textbf{1} & \textbf{2} & \textbf{3} & \textbf{4} & \textbf{5} \\ \hline
\textbf{الإجابة} &  &  &  &  &  \\[8pt] \hline
\end{tabular}

\vspace{8pt}
\renewcommand{\arraystretch}{2.2}
\begin{tabular}{|>{\centering\arraybackslash}m{7.5cm}|>{\centering\arraybackslash}m{3.2cm}|>{\centering\arraybackslash}m{3.2cm}|}
\hline
\textbf{الاسم واللقب} & \textbf{الفوج} & \textbf{العلامة} \\ \hline
 & & \hfill /20 \\ \hline
\end{tabular}
\end{center}
\end{document}
